Os Large Language Models (LLMs) têm se consolidado como uma ferramenta central na obtenção e interpretação de informações, exercendo impactos significativos no cotidiano social. Já se observam aplicações relevantes desses modelos em áreas como a Educação \cite{IAeChatGPT}, com avanços nos métodos de ensino e o uso como ferramenta de auxilio pedagógico, no campo da medicina \cite{ChatGPTInMedicice}, oferecendo suporte à análise de dados clinicos e à consulta de literatura na biomédica, e na economia \cite{EconomieAndLLM}, com a realização de análises do mercado de ações e avaliações qualitativas de empresas para fins de investimento.


Na pecuária leiteira já é desejado que um produtor agropecuário possa fazer perguntas ao modelo e seja retornado respostas coerentes e fundamentadas sobre a pecuária leiteira. Uma aplicação prática relevante seria a geração de respostas baseadas em literatura científica para questões relacionadas à alimentação do rebanho e ao aumento da produção de leite. Para tal, o LLM deve ser imbuido com conhecimento especializado, utilizando textos científicos ou conjuntos de dados relacionadas à lactação e ao manejo dos animais.


Nesse contexto, técnicas como a Geração Aumentada por Recuperação (RAG) \cite{RAG} tornam-se essenciais, pois permitem adaptar modelos existentes a um domínio em específico, incorporando conhecimento especializado que pode não estar presente no modelo original. A avaliação pode ser realizada por meio de métricas de avaliação textual consolidadas, como BLEU \cite{papineni-etal-2002-bleu}, ROUGE \cite{lin-2004-rouge}, BERTScore \cite{zhang2020bertscoreevaluatingtextgeneration}, BARTScore \cite{bartscore} e GPTScore \cite{fu2023gptscoreevaluatedesire}, as quais possibilitam medir a similaridade entre as respostas geradas pelos modelos e os textos de referência.


Desse modo, torna-se necessária a análise de aplicação de LLMs associadas a uma base de conhecimento específica, considerando o potencial como ferramenta de grande auxílio de suporte aos pecuaristas. Essa abordagem pode viabilizar a geração de respostas rápidas, detalhadas e adaptadas de forma aprofundada ao contexto da produção leiteira. 